% Figure — Référence de tension MOS à double tension de seuil (dual-V_th) : miroir PMOS, M1 monté en diode, M2 dégénéré par R_Δ (EANA3, chapitre 5).
% D'après les notes externes [notes-2023-c10, p. 5] et le corrigé externe [corr-2023-ref-cmos, p. 1-2] (notations R et R_Δ), redessinée.
% Rendu SVG :
%   latex -output-directory=<dossier temporaire> 05-reference-dual-vth.tex
%   dvisvgm --no-fonts --exact-bbox -o 05-reference-dual-vth.svg <dossier temporaire>/05-reference-dual-vth.dvi
\documentclass[border=4pt]{standalone}
\usepackage[T1]{fontenc}
\usepackage[utf8]{inputenc}
\usepackage{amsmath}
\usepackage{sansmath}
\renewcommand{\familydefault}{\sfdefault}
\usepackage[european, straightvoltages, americancurrents]{circuitikz}
\usetikzlibrary{shapes.geometric}
\begin{document}
\sansmath
\begin{circuitikz}[font=\small, >=latex, fleche/.style={->, thick}]
  \draw (0,5) node[pmos, xscale=-1] (Ma) {};  \node[left] at (-0.45,5) {$M_{101}$};
  \draw (3,5) node[pmos] (Mb) {};  \node[right] at (3.45,5) {$M_{102}$};
  \draw (-0.4,6.1) -- (3.4,6.1);  \node[above] at (1.5,6.1) {$V_{DD}$};
  \draw (Ma.S) -- (0,6.1);  \draw (Mb.S) -- (3,6.1);
  \draw (Ma.G) -- (Mb.G);
  \draw (1.9,5) node[circ]{} -- (1.9,3.8) -- (3,3.8) node[circ]{};
  \draw (Ma.D) -- (0,3.6) coordinate (vr);
  \draw (vr) node[circ]{} to[R, l_=$R$, i>^=$I$] (0,2.0) coordinate (d1);
  \draw (vr) -- (-1.3,3.6) node[ocirc]{} node[left]{$V_{ref}$};
  \draw (0,1) node[nmos, xscale=-1] (M1) {};  \node[left] at (-0.45,1) {$M_1$};
  \draw (d1) node[circ]{} -- (M1.D);
  \draw (M1.G) -- (1.3,1) node[circ]{} -- (1.3,2.0) -- (d1);
  \draw (3,1) node[nmos] (M2) {};  \node[right] at (3.45,1) {$M_2$};
  \draw (1.3,1) -- (M2.G);
  \draw (Mb.D) -- (M2.D);
  \draw (M2.S) to[R, l=$R_\Delta$] (3,-1.3);
  \draw (M1.S) -- (0,-1.3);  \draw (-1.3,-1.3) node[ocirc]{} -- (3,-1.3);  \draw (1.5,-1.3) node[ground]{};
  \draw[fleche] (3.9,-1.1) -- node[right]{$\Delta V_{th}$} (3.9,0.1);
\end{circuitikz}
\end{document}
